
\section{Contribuciones propuestas y estatus de novedad}

Las siguientes piezas deben entenderse como candidatas de contribución del programa IRL. No se afirma todavía novedad matemática cerrada para ninguna de ellas sin una revisión sistemática contra inverse procedural modeling, 4D BIM, assembly/disassembly planning, Harris matrices, computational archaeology, Bayesian chronology y multiparameter persistence.

\begin{enumerate}[label=\textbf{C\arabic*.},leftmargin=*]
\item \textbf{Inversión terminal multimodal}: integración de estado terminal, orden parcial, física, accesibilidad y evidencia en una distribución de historias.
\item \textbf{Ontología dual}: \(\mathcal A=(\Kp,\Kn)\), con materia y espacio negativo como entidades de igual rango ontológico.
\item \textbf{Masa ausente contrafactual}: \(m^\ominus\) como interfaz entre vacío arquitectónico y densidad de referencia.
\item \textbf{Campos duales de nacimiento}: \(T_b^+\) y \(T_b^-\).
\item \textbf{Bifiltración negativa}: \((\tau,r)\mapsto V_{\tau,r}\).
\item \textbf{Frente inverso admisible}: máximos de un poset intersectados con mecánica, acceso y evidencia.
\item \textbf{Ontología temporal latente}: inferencia de clases funcionales de objetos desaparecidos.
\item \textbf{Jerarquía I0--I3}: resolución explícita de identificabilidad histórica.
\item \textbf{Protocolo residual}: cálculo de \(\Fres\) sólo después de gates de identificabilidad y agotamiento de procesos convencionales.
\item \textbf{Triada directo--inverso--contrafactual}: toda historia debe volver a producir observaciones potencialmente falsables.
\end{enumerate}

\section{Conjeturas de investigación}

\begin{conjecture}[Contracción multimodal]
Bajo likelihoods correctamente especificados y observaciones consistentes, nuevas modalidades independientes deberían reducir el volumen efectivo de historias admisibles, salvo que revelen misspecification del modelo.
\end{conjecture}

\begin{conjecture}[Ventaja dual]
Para arquitecturas con vacíos internos significativos, la inferencia conjunta sobre \((\Kp,\Kn)\) discrimina cronologías mejor que la geometría material aislada.
\end{conjecture}

\begin{conjecture}[Valor de la ontología temporal]
Permitir \(\mathcal T\) reduce falsos residuos de fuerza respecto de un solver limitado al estado terminal material.
\end{conjecture}

\begin{conjecture}[Diseño activo]
Seleccionar mediciones mediante EIG discrimina historias supervivientes con menor costo experimental que aumentar uniformemente la resolución de todas las modalidades.
\end{conjecture}

\begin{conjecture}[Residuo físico robusto]
Sólo un residuo que permanezca significativamente distinto de cero bajo refinamiento ontológico, mecanismos convencionales, variación razonable de parámetros, propagación de incertidumbre y análisis de identificabilidad merece elevarse a hipótesis de interacción no modelada.
\end{conjecture}

\section{Programa experimental para Khufu}

El primer experimento serio no debería intentar resolver la pirámide completa. Se propone comenzar por el subsistema de la Cámara del Rey:

\begin{enumerate}
\item construir un modelo 3D documentado de cámara, acceso, techo, cámaras de descarga, estructura superior y masonry circundante;
\item clasificar cada entidad como observada, derivada, inferida, hipótesis o desconocida;
\item construir \(\Kp\) y \(\Kn\);
\item estimar \(T_b^+\) y \(T_b^-\) como campos latentes;
\item construir un grafo tipado de dependencias;
\item enumerar o muestrear extensiones lineales;
\item aplicar mecánica de contacto y estabilidad incremental;
\item evaluar accesibilidad de piezas en \(\SE(3)\);
\item calcular eventos topológicos de \(V_{\tau,r}\);
\item integrar procedencia y rutas cuando exista evidencia;
\item generar forward predictions de densidad y muografía;
\item identificar equivalencias de historias;
\item usar EIG para priorizar nuevas observaciones.
\end{enumerate}

El outcome principal debe ser cuantitativo:

\begin{equation}
\boxed{
\text{¿qué fracción del espacio de historias elimina cada gate independiente?}
}
\end{equation}

\section{Matriz de suficiencia observacional}

Antes de calcular una fuerza residual debe publicarse una matriz de modalidades disponibles:

\begin{equation}
\mathcal O_{\mathrm{available}}
=
[
G,
\text{microestructura},
M,
L,
S,
C,
F,
E,
H,
Q,
A,
T,
\ldots
].
\end{equation}

Cada cantidad o claim se etiqueta como:

\begin{description}
\item[OBSERVED:] medida directa;
\item[DERIVED:] calculada desde observaciones;
\item[INFERRED:] posterior dependiente del modelo;
\item[HYPOTHESIS:] posibilidad no discriminada;
\item[COUNTERFACTUAL:] resultado de una intervención simulada;
\item[UNKNOWN:] no restringido;
\item[NON-IDENTIFIABLE:] múltiples estados producen evidencia indistinguible.
\end{description}

\section{Amenazas a la validez}

\subsection{Sesgo retrospectivo}

Los cinco benchmarks se diseñaron con conocimiento general de las estructuras. Aunque se separó entrada terminal y revelación histórica, no constituyen experimentos doble ciego.

\subsection{Granularidad manual}

Los grafos son coarse-grained. Resultados combinatorios como \(99.999995\%\) describen reducción dentro de un poset seleccionado, no precisión histórica.

\subsection{Priors de ingeniería}

Ciertas inferencias reflejan conocimientos de ingeniería incorporados en reglas. Deben externalizarse, versionarse y compararse con baselines.

\subsection{Equifinalidad}

Múltiples mecanismos pueden producir el mismo terminal. El formalismo debe preservar esa pluralidad.

\subsection{Calidad de fuentes}

Resúmenes institucionales pueden contener simplificaciones o discrepancias. El caso del parámetro métrico de la esfera de Sydney llevó a excluir ese número de la puntuación hasta disponer de una fuente geométrica primaria.

\subsection{Flexibilidad del campo residual}

Un campo hipotético muy flexible puede ajustar casi cualquier historia. Debe recibir fuerte penalización de complejidad y generar predicciones independientes.

\section{Criterios de falsificación del programa}

IRL debe considerarse debilitado si ocurre sistemáticamente alguno de los siguientes resultados:

\begin{enumerate}
\item no discrimina historias mejor que baselines simples de altura o soporte;
\item introduce numerosas entidades temporales ad hoc sin mejora predictiva fuera de muestra;
\item falla en recuperar relaciones causales gruesas en benchmarks preregistrados;
\item sus posteriors no calibran incertidumbre;
\item no predice observaciones ocultas mejor que alternativas convencionales;
\item un supuesto residuo físico desaparece al incorporar una modalidad omitida;
\item los formalismos propuestos resultan equivalentes a herramientas existentes sin aportar capacidad inferencial.
\end{enumerate}

\section{Plan de validación prospectiva}

La fase siguiente debe cambiar de retrospectiva a prospectiva:

\begin{enumerate}
\item seleccionar estructuras documentadas adicionales;
\item congelar reglas, ontología y scoring;
\item entregar al solver sólo datasets terminales definidos de antemano;
\item ocultar cronología y maquinaria;
\item registrar predicciones antes de revelar ground truth;
\item medir precision/recall de precedencias;
\item medir calibración de clases funcionales;
\item penalizar complejidad de \(\mathcal T\);
\item comparar con baselines de altura, soporte y planners convencionales;
\item publicar fallos, no sólo éxitos.
\end{enumerate}

Sólo después tiene sentido interpretar resultados sobre Khufu como algo más que una demostración conceptual.
